\documentclass[12pt,a4paper]{report}
\input{common/volume-preamble}

\begin{document}

\chapter*{Model Box Sample}

\begin{modelbox}{Model (Classical Propositional Logic)}
\begin{model}[Classical Propositional Logic]\label{model:classical-propositional-logic}
\modeldepends{
  def:propositional-language,
  def:truth-value-propositional-logic,
  def:truth-assignment-propositional-logic,
  def:satisfaction-propositional-formula
}

Classical propositional logic is presented as a two-valued model whose
syntax is given by a propositional language and whose semantics is given by
valuations into the two-element Boolean algebra.

\begin{modelsorts}
  \modelsort{\mathcal P}{\text{a set of propositional variables}}
    {The syntactic atoms of the language.}
  \modelsort{\Phi}{\text{the set of well-formed formulas over }\mathcal P}
    {The formulas generated from atoms and connectives.}
  \modelsort{\mathbf 2}{\{0,1\}}
    {The semantic truth-value domain.}
\end{modelsorts}

\begin{modelconstants}
  \modelconstant{0}{def:truth-value-propositional-logic}
  \modelconstant{1}{def:truth-value-propositional-logic}
\end{modelconstants}

\begin{modelarity}
  \arity{\neg_{\Phi}}{1}
  \arity{\land_{\Phi}}{2}
  \arity{\lor_{\Phi}}{2}
  \arity{\to_{\Phi}}{2}
  \arity{\leftrightarrow_{\Phi}}{2}
\end{modelarity}

\begin{modelfunctions}
  \modelfunction{\neg^{\mathbf 2}}{\mathbf 2 \to \mathbf 2}
    {Truth-value negation.}
  \modelfunction{\land^{\mathbf 2}}{\mathbf 2^2 \to \mathbf 2}
    {Truth-value conjunction.}
  \modelfunction{\lor^{\mathbf 2}}{\mathbf 2^2 \to \mathbf 2}
    {Truth-value disjunction.}
  \modelfunction{\to^{\mathbf 2}}{\mathbf 2^2 \to \mathbf 2}
    {Truth-value implication.}
  \modelfunction{\leftrightarrow^{\mathbf 2}}{\mathbf 2^2 \to \mathbf 2}
    {Truth-value biconditional.}
\end{modelfunctions}

\begin{modelrelations}
  \modelrelation{\models}{v \models \varphi}
    {Satisfaction of a formula by a valuation.}
  \modelrelation{\models}{\Gamma \models \varphi}
    {Semantic consequence from a set of assumptions.}
\end{modelrelations}

\begin{modelsemantics}
  \modelclause{A valuation is a function $v:\mathcal P\to\mathbf 2$.}
  \modelclause{Satisfaction is defined by $v\models\varphi$ iff
  $\bar v(\varphi)=1$.}
  \modelclause{Semantic consequence is defined by $\Gamma\models\varphi$ iff
  for every valuation $v:\mathcal P\to\mathbf 2$, if $v\models\gamma$ for
  every $\gamma\in\Gamma$, then $v\models\varphi$.}
\end{modelsemantics}

\begin{modelbridge}
  \bridgeclause{$\bar v(p)=v(p)$ for every atomic variable $p\in\mathcal P$.}
  \bridgeclause{$\bar v(\neg_{\Phi}\varphi)=
  \neg^{\mathbf 2}\bar v(\varphi)$.}
  \bridgeclause{$\bar v(\varphi\land_{\Phi}\psi)=
  \land^{\mathbf 2}(\bar v(\varphi),\bar v(\psi))$.}
  \bridgeclause{$\bar v(\varphi\lor_{\Phi}\psi)=
  \lor^{\mathbf 2}(\bar v(\varphi),\bar v(\psi))$.}
  \bridgeclause{$\bar v(\varphi\to_{\Phi}\psi)=
  \to^{\mathbf 2}(\bar v(\varphi),\bar v(\psi))$.}
  \bridgeclause{$\bar v(\varphi\leftrightarrow_{\Phi}\psi)=
  \leftrightarrow^{\mathbf 2}(\bar v(\varphi),\bar v(\psi))$.}
\end{modelbridge}

\begin{modelaxioms}
  \modelaxiom{The semantic algebra is the two-element Boolean algebra
  $\mathbf 2=(\{0,1\},\neg^{\mathbf 2},\land^{\mathbf 2},
  \lor^{\mathbf 2},\to^{\mathbf 2},\leftrightarrow^{\mathbf 2})$.}
\end{modelaxioms}

\begin{modelconsequences}
  \modelconsequence{Every valuation $v:\mathcal P\to\mathbf 2$ extends
  uniquely to a function $\bar v:\Phi\to\mathbf 2$ satisfying the bridge
  clauses.}
\end{modelconsequences}
\end{model}
\end{modelbox}

\end{document}
